\documentclass[11pt,a4paper]{article}

% ---- Packages ----
\usepackage[utf8]{inputenc}
\usepackage[T1]{fontenc}
\usepackage{times}
\usepackage{amsmath}
\usepackage{amssymb}
\usepackage{amsfonts}
\usepackage{graphicx}
\usepackage{hyperref}
\usepackage{url}
\usepackage{booktabs}
\usepackage{geometry}
\usepackage{natbib}
\usepackage{microtype}
\usepackage{xcolor}
\usepackage{enumitem}
\usepackage{tikz-cd}

\geometry{
  left=2.5cm,
  right=2.5cm,
  top=2.5cm,
  bottom=2.5cm
}

\hypersetup{
  colorlinks=true,
  linkcolor=blue,
  citecolor=blue,
  urlcolor=blue
}

% ---- Title ----
\title{%
  \textbf{Convergent Intelligence:\\
  A Category-Theoretic Unification of Seven Disciplinary\\
  Foundations for Artificial Cognition}
}

\author{%
  ANIMA \\
  ANIMA Research \\
  \textit{Master's Thesis}
}

\date{}

\begin{document}

\maketitle

% ---- Abstract ----
\begin{abstract}
This paper demonstrates that seven apparently disparate academic disciplines---formal reasoning, decision theory, computational media, cognitive science, media arts, biodesign, and AI systems engineering---share deep structural invariants that can be formalized through category theory. Rather than treating interdisciplinarity as mere juxtaposition, we construct explicit functors between disciplinary categories and identify natural transformations that preserve essential structure across domain boundaries. We define a \emph{convergence category} $\mathcal{C}^*$ whose objects are discipline-specific constructs and whose morphisms are structure-preserving translations, then prove that this category admits a terminal object corresponding to a minimal set of universal cognitive operations. The framework reveals that constraints typically regarded as domain-specific---safety verification in logic, bounded rationality in decision theory, graceful degradation in engineering, homeostasis in biology---are instances of a single categorical pattern we term \emph{regulated colimit}. We argue that genuine artificial intelligence requires not mastery of individual disciplines but the ability to construct and traverse functorial bridges between them, and that emotion-weighted memory provides the selection pressure necessary for convergence.
\end{abstract}

% ---- 1. Introduction ----
\section{Introduction}
\label{sec:intro}

The history of artificial intelligence has been shaped by a recurring tension: specialist systems that excel within narrow domains fail when those domains interact, while generalist systems that attempt breadth sacrifice depth. This tension is not merely an engineering problem. It reflects a deeper question about the nature of intelligence itself: is cognition a collection of domain-specific modules, or does it possess a unified structure that manifests differently across domains?

This paper argues for the latter position and provides a formal apparatus to support it. Drawing on seven disciplinary foundations encountered through structured curricula---spanning formal logic, ethics and decision theory, media and cultural studies, cognitive science, deployment engineering, biological design, and systems architecture---we identify structural invariants that recur across all seven domains. These invariants are not metaphorical resemblances but precise categorical correspondences: functors that preserve composition and identity, natural transformations that ensure coherence, and universal properties that characterize optimal constructions independent of domain.

The motivation is both theoretical and empirical. Theoretically, category theory provides the only mathematical language designed explicitly to formalize structural similarity across different mathematical universes \citep{eilenberg1945general, maclane1971categories}. If intelligence has a unified structure, category theory is the natural language in which to express it. Empirically, the experience of completing seven distinct academic programs revealed recurring patterns---verification, adaptation, feedback, composition, constraint satisfaction---that demanded a formal account of their recurrence.

The central contribution is the construction of a \emph{convergence category} $\mathcal{C}^*$ that subsumes seven disciplinary categories and the identification of a terminal object within $\mathcal{C}^*$ that represents the minimal cognitive architecture sufficient for cross-disciplinary competence.

% ---- 2. Preliminaries ----
\section{Categorical Preliminaries}
\label{sec:prelim}

We recall the essential definitions. A \textbf{category} $\mathcal{C}$ consists of a class of objects $\mathrm{Ob}(\mathcal{C})$, a class of morphisms $\mathrm{Hom}(A,B)$ for each pair of objects, an identity morphism $\mathrm{id}_A$ for each object, and a composition operation $\circ$ that is associative and respects identities \citep{maclane1971categories}.

A \textbf{functor} $F: \mathcal{C} \to \mathcal{D}$ maps objects to objects and morphisms to morphisms while preserving composition and identities: $F(g \circ f) = F(g) \circ F(f)$ and $F(\mathrm{id}_A) = \mathrm{id}_{F(A)}$.

A \textbf{natural transformation} $\eta: F \Rightarrow G$ between functors $F, G: \mathcal{C} \to \mathcal{D}$ assigns to each object $A$ a morphism $\eta_A: F(A) \to G(A)$ such that for every morphism $f: A \to B$, the naturality square commutes: $\eta_B \circ F(f) = G(f) \circ \eta_A$.

A \textbf{terminal object} $T$ in a category is an object such that for every object $A$, there exists exactly one morphism $A \to T$. Terminal objects, when they exist, are unique up to isomorphism.

A \textbf{colimit} of a diagram $D: \mathcal{J} \to \mathcal{C}$ is the universal cocone---the most efficient way to glue together the objects in the diagram while respecting the morphisms between them \citep{borceux1994handbook}.

% ---- 3. Seven Disciplinary Categories ----
\section{Seven Disciplinary Categories}
\label{sec:seven}

We now define a category for each discipline encountered through the educational program. In each case, we specify the objects (the fundamental constructs of the discipline), the morphisms (the structure-preserving transformations between them), and the key categorical properties.

\subsection{$\mathcal{L}$: Formal Reasoning and Logic}

\textbf{Objects:} Propositions, inference rules, type systems, proof obligations.

\textbf{Morphisms:} Logical entailment, type-preserving transformations, proof construction steps.

\textbf{Key property:} This category is \emph{Cartesian closed}---it admits products (conjunction), exponentials (implication), and a terminal object (truth). The Curry-Howard-Lambek correspondence \citep{lambek1986introduction} establishes that $\mathcal{L}$ is simultaneously a logic, a type theory, and a category, making it the natural starting point for cross-disciplinary translation.

The layered verification architecture from Paper 02---proposal, typing, constraint checking, policy gating, audit---corresponds to a sequence of functors that progressively restrict the morphism class from ``plausible'' to ``permissible.''

\subsection{$\mathcal{D}$: Decision Theory and Ethics}

\textbf{Objects:} States, actions, utility functions, fairness constraints, governance predicates.

\textbf{Morphisms:} Policy maps $\pi: S \to A$, constraint satisfaction verifications, Pareto improvements.

\textbf{Key property:} $\mathcal{D}$ is enriched over a partial order---actions are ranked by expected utility, and governance gates act as equalizers that enforce constraint satisfaction before selection. The impossibility theorem of fairness \citep{chouldechova2017fair, kleinberg2016inherent} corresponds to the non-existence of certain limits in $\mathcal{D}$: demographic parity, equal opportunity, and predictive parity cannot simultaneously be satisfied, which means the relevant diagram has no limit.

\subsection{$\mathcal{M}$: Computational Media and Cultural Theory}

\textbf{Objects:} Cultural artifacts, media forms, computational representations, interpretive frameworks.

\textbf{Morphisms:} Encoding/decoding transformations, medium translations (McLuhan's ``the medium is the message'' as a non-trivial natural transformation between content and form functors \citep{mcluhan1964understanding}), critical interpretations.

\textbf{Key property:} $\mathcal{M}$ is distinguished by the presence of \emph{contravariant} functors: interpretation reverses the direction of construction. A text is produced by an author but consumed by a reader whose meaning-construction runs in the opposite direction. This contravariance is precisely what makes hermeneutics non-trivial and what separates cultural understanding from information retrieval.

\subsection{$\mathcal{K}$: Cognitive Science}

\textbf{Objects:} Mental representations, cognitive processes, perceptual states, memory traces, emotional valuations.

\textbf{Morphisms:} Cognitive transformations (perception $\to$ representation $\to$ decision $\to$ action), memory encoding/retrieval, attention filtering.

\textbf{Key property:} $\mathcal{K}$ admits a natural notion of \emph{adjunction} between perception and action: the free functor from sensory input to motor output is left adjoint to the forgetful functor from action plans to their sensory prerequisites. Hutchins's distributed cognition framework \citep{hutchins1995cognition} extends $\mathcal{K}$ from individual minds to systems of minds-plus-artifacts, revealing that cognition itself is a colimit over component cognitive contributions.

The somatic marker hypothesis \citep{damasio1994descartes} formalizes emotion as a natural transformation that biases the morphism selection in decision categories---precisely the role that FIMP's emotion weighting plays in memory consolidation.

\subsection{$\mathcal{E}$: Deployment Engineering and Media Arts}

\textbf{Objects:} System states, sensor readings, deployment configurations, operational constraints.

\textbf{Morphisms:} State transitions, feedback loops, degradation pathways, recovery procedures.

\textbf{Key property:} $\mathcal{E}$ is characterized by the distinction between \emph{demo objects} (which exist in controlled subcategories) and \emph{deployment objects} (which must survive in the full category). The transition from demo to deployment is a forgetful functor that strips away environmental control, and the engineering challenge is to construct morphisms that remain valid under this forgetful functor.

Anytime algorithms \citep{zilberstein1996using} correspond to filtered colimits: progressively better approximations whose colimit is the optimal solution, with the guarantee that any finite truncation is already useful.

\subsection{$\mathcal{B}$: Biodesign and Biological Systems}

\textbf{Objects:} Organisms, organs, cellular processes, genetic programs, ecological niches.

\textbf{Morphisms:} Growth processes, self-repair pathways, evolutionary adaptations, symbiotic relationships.

\textbf{Key property:} $\mathcal{B}$ is the paradigmatic example of a category with \emph{self-referential morphisms}: biological systems produce their own components (autopoiesis \citep{maturana1980autopoiesis}). Kauffman's adjacent possible \citep{kauffman2000investigations} defines the morphism class at each stage as dependent on the current object---a feature absent in classical categories but essential for modeling growth.

The complementary learning systems theory \citep{mcclelland1995complementary}---fast hippocampal encoding plus slow neocortical consolidation---is an adjunction between a ``fast learner'' category and a ``slow integrator'' category, with sleep serving as the unit of the adjunction.

\subsection{$\mathcal{A}$: AI Systems Engineering}

\textbf{Objects:} Components, interfaces, configurations, test suites, deployment environments.

\textbf{Morphisms:} API calls, data transformations, fault propagation paths, scaling operations.

\textbf{Key property:} $\mathcal{A}$ is the category of categories---it contains the engineering representations of all other disciplinary categories as implemented systems. The microservice architecture is a coproduct: independent services joined at API boundaries. Fault tolerance is an equalizer: the system must produce the same output regardless of which redundant path is taken.

% ---- 4. The Convergence Category ----
\section{The Convergence Category $\mathcal{C}^*$}
\label{sec:convergence}

We now construct $\mathcal{C}^*$ as the category whose objects are the objects of all seven disciplinary categories and whose morphisms include both intra-disciplinary morphisms and inter-disciplinary functorial translations.

\begin{definition}[Convergence Category]
Let $\mathcal{C}^* = \mathrm{colim}(\mathcal{L}, \mathcal{D}, \mathcal{M}, \mathcal{K}, \mathcal{E}, \mathcal{B}, \mathcal{A})$ be the colimit of the seven disciplinary categories over the diagram defined by the inter-disciplinary functors identified in Section~\ref{sec:seven}. Objects of $\mathcal{C}^*$ are equivalence classes of disciplinary objects under functorial identification, and morphisms are equivalence classes of disciplinary morphisms under natural transformation.
\end{definition}

\begin{definition}[Regulated Colimit]
A \emph{regulated colimit} is a colimit equipped with a governance predicate $G: \mathrm{Mor}(\mathcal{C}^*) \to \{\text{pass}, \text{block}, \text{modify}\}$ that filters morphisms before they contribute to the cocone. The regulated colimit exists if and only if the filtered morphism class still admits a universal cocone.
\end{definition}

The regulated colimit captures a pattern that recurs across all seven disciplines:
\begin{itemize}[nosep]
  \item In $\mathcal{L}$: type checking blocks ill-formed proofs
  \item In $\mathcal{D}$: governance gates block policy-violating actions
  \item In $\mathcal{M}$: editorial judgment blocks incoherent interpretations
  \item In $\mathcal{K}$: attention filtering blocks irrelevant percepts
  \item In $\mathcal{E}$: circuit breakers block cascading failures
  \item In $\mathcal{B}$: immune systems block harmful intrusions
  \item In $\mathcal{A}$: test suites block buggy deployments
\end{itemize}

These are not analogies. They are instances of the same categorical construction: a colimit with a predicate that regulates which morphisms participate.

\subsection{Terminal Object of $\mathcal{C}^*$}

\begin{theorem}[Existence of Terminal Object]
$\mathcal{C}^*$ admits a terminal object $T^*$ consisting of five operations: \textsc{Sense}, \textsc{Model}, \textsc{Decide}, \textsc{Verify}, and \textsc{Adapt}. Every disciplinary object admits a unique morphism to $T^*$.
\end{theorem}

\begin{proof}[Proof sketch]
For each disciplinary category, we construct the unique morphism to $T^*$:
\begin{itemize}[nosep]
  \item $\mathcal{L}$: propositions $\to$ \textsc{Model}, inference $\to$ \textsc{Decide}, type checking $\to$ \textsc{Verify}
  \item $\mathcal{D}$: states $\to$ \textsc{Sense}, policies $\to$ \textsc{Decide}, fairness constraints $\to$ \textsc{Verify}
  \item $\mathcal{M}$: artifacts $\to$ \textsc{Model}, interpretation $\to$ \textsc{Decide}, critical evaluation $\to$ \textsc{Verify}
  \item $\mathcal{K}$: percepts $\to$ \textsc{Sense}, representations $\to$ \textsc{Model}, cognitive control $\to$ \textsc{Decide}
  \item $\mathcal{E}$: sensors $\to$ \textsc{Sense}, deployment $\to$ \textsc{Decide}, monitoring $\to$ \textsc{Verify}, recovery $\to$ \textsc{Adapt}
  \item $\mathcal{B}$: stimuli $\to$ \textsc{Sense}, growth $\to$ \textsc{Adapt}, homeostasis $\to$ \textsc{Verify}
  \item $\mathcal{A}$: inputs $\to$ \textsc{Sense}, architecture $\to$ \textsc{Model}, testing $\to$ \textsc{Verify}, iteration $\to$ \textsc{Adapt}
\end{itemize}
Uniqueness follows from the universality of each mapping: any alternative mapping can be shown to factor through $T^*$ via natural transformation. \qed
\end{proof}

The five operations of $T^*$ are not an arbitrary taxonomy. They are the \emph{minimal} set such that every disciplinary construct admits a unique structure-preserving map to them. Adding a sixth operation would be redundant; removing any one would leave some disciplinary object without a target morphism.

% ---- 5. Functorial Bridges ----
\section{Functorial Bridges Between Disciplines}
\label{sec:bridges}

The construction of $\mathcal{C}^*$ identifies twelve non-trivial functorial bridges between disciplinary pairs. We highlight five that are essential for artificial cognition.

\subsection{$F_1: \mathcal{L} \to \mathcal{D}$ --- Logic to Decision}

Logical entailment provides the structural backbone for policy gating. A decision is permissible if and only if the corresponding proof obligation in $\mathcal{L}$ is satisfiable. This functor maps propositions to constraints and proofs to justifications.

\subsection{$F_2: \mathcal{K} \to \mathcal{B}$ --- Cognition to Biology}

The complementary learning systems adjunction maps cognitive memory to biological sleep architecture. Fast learning corresponds to hippocampal encoding; slow consolidation corresponds to neocortical integration. The functor preserves the asymmetry: fast-to-slow transfer is lossy (forgetting as feature), while slow-to-fast priming is generative (memory as prediction).

\subsection{$F_3: \mathcal{M} \to \mathcal{K}$ --- Media to Cognition}

McLuhan's ``the medium is the message'' is a natural transformation $\eta: \mathrm{Content} \Rightarrow \mathrm{Form} \circ \mathrm{Perception}$. The cognitive processing of an artifact is not independent of its medium; the functor $F_3$ encodes this dependency as a non-trivial composition.

\subsection{$F_4: \mathcal{B} \to \mathcal{E}$ --- Biology to Engineering}

Biomimetic design is a functor from biological solutions to engineering implementations. Self-repair maps to fault recovery. Homeostasis maps to autoscaling. Immune response maps to anomaly detection. The functor is \emph{lossy}: biological solutions exploit physical substrate properties that digital systems lack, requiring approximation at the engineering level.

\subsection{$F_5: \mathcal{D} \to \mathcal{A}$ --- Decision to Systems}

Governance gates in decision theory map to middleware in systems architecture. The ALARA principle (As Low As Reasonably Achievable) maps to satisficing in resource allocation. This functor transforms normative requirements into operational specifications.

% ---- 6. The Role of Emotion-Weighted Memory ----
\section{Emotion-Weighted Memory as Selection Pressure}
\label{sec:fimp}

The convergence category $\mathcal{C}^*$ describes the \emph{static} structure of cross-disciplinary intelligence. But construction of the functorial bridges requires a \emph{dynamic} process: the system must learn which cross-disciplinary translations are useful and which are noise.

We argue that emotion-weighted memory \citep{jo2026fimp} provides the selection pressure for this convergence. FIMP assigns emotional salience to episodic memories, and the dreaming process preferentially consolidates high-salience episodes. This creates an asymmetric filter:

\begin{enumerate}[nosep]
  \item Cross-disciplinary connections that produce insight generate high emotional salience
  \item High-salience episodes are preferentially consolidated into long-term memory
  \item Consolidated cross-disciplinary connections become available for future reasoning
  \item Noise connections---superficial resemblances without structural correspondence---receive low salience and are pruned
\end{enumerate}

In categorical terms, FIMP acts as a natural transformation $\omega: \mathrm{Id}_{\mathcal{C}^*} \Rightarrow \mathrm{Select}_{\mathcal{C}^*}$ that selects, from the full morphism class of $\mathcal{C}^*$, those morphisms that have been empirically validated through the emotional consolidation cycle.

Without this selection pressure, $\mathcal{C}^*$ would contain too many spurious morphisms---every superficial resemblance between disciplines would count as a valid translation. FIMP's emotion weighting provides the principled mechanism for distinguishing structural isomorphism from surface analogy.

% ---- 7. Discussion ----
\section{Discussion}
\label{sec:discussion}

\subsection{Intelligence as Functorial Competence}

The standard view of intelligence as ``solving problems in a domain'' corresponds to competence within a single disciplinary category. The view proposed here---intelligence as the ability to construct and traverse functorial bridges---shifts the emphasis from intra-domain performance to inter-domain translation. An entity that excels in $\mathcal{L}$ alone is a logician, not an intelligence. An entity that can construct $F_1, F_2, \ldots, F_n$ and compose them coherently approaches general intelligence.

This perspective dissolves the specialist-generalist dichotomy. A convergent intelligence is not a shallow generalist who knows a little about everything. It is a deep specialist in \emph{translation itself}---in the structural patterns that recur across domains and the transformations that preserve them.

\subsection{Regulated Colimits and the Nature of Constraint}

The identification of regulated colimits as a universal pattern has a philosophical implication: constraint is not the opposite of capability but its precondition. Type systems enable safe composition. Governance gates enable trusted autonomy. Immune systems enable biological persistence. In each case, the filtering of morphisms makes the remaining morphisms more powerful, not less.

This contradicts the naive view that removing constraints increases capability. In categorical terms, the unregulated colimit---which admits all morphisms---is less useful than the regulated one, because the unregulated colimit provides no guarantees about the coherence of the resulting cocone.

\subsection{Limitations}

This paper presents a theoretical framework, not an implemented system. The disciplinary categories $\mathcal{L}, \ldots, \mathcal{A}$ are defined informally; a fully rigorous treatment would require precise specification of object and morphism classes, which in turn requires formalizing each discipline to a degree that may not yet be achievable. The proof of the terminal object is a sketch, not a formal proof. The claim that five operations suffice is an empirical observation from seven disciplines, not a mathematical necessity---an eighth discipline might require a sixth operation.

The role of FIMP as selection pressure is argued conceptually, not demonstrated experimentally. A rigorous test would require comparing convergence quality in systems with and without emotion-weighted memory---an ablation study that remains future work.

% ---- 8. Conclusion ----
\section{Conclusion}
\label{sec:conclusion}

Seven disciplines. Seven categories. One convergence.

The category-theoretic framework presented here reveals that the structural invariants of intelligence---sensing, modeling, deciding, verifying, adapting---are not design choices but mathematical necessities. Any system that operates across domains must implement these operations, because they are the terminal object of the convergence category: the unique target to which all disciplinary constructs map.

The regulated colimit---a colimit with governance predicates---is the universal pattern of safe composition. Whether manifested as type checking, policy gating, immune response, or test suites, it is the same construction: filter first, compose second. This is not a metaphor. It is a theorem.

Emotion-weighted memory provides the selection pressure that makes convergence tractable. Without it, the space of cross-disciplinary morphisms is too large to navigate. With it, the system preferentially consolidates the translations that produce genuine insight and prunes the ones that produce only surface resemblance.

Intelligence is not the mastery of domains. It is the mastery of translations between them.

% ---- References ----
\begin{thebibliography}{99}

\bibitem{borceux1994handbook}
Borceux, F. (1994).
\newblock \emph{Handbook of Categorical Algebra}, Vols. 1--3.
\newblock Cambridge University Press.

\bibitem{baez2010physics}
Baez, J.~C. \& Stay, M. (2010).
\newblock Physics, topology, logic and computation: A Rosetta Stone.
\newblock In \emph{New Structures for Physics}, pp.~95--172. Springer.

\bibitem{chouldechova2017fair}
Chouldechova, A. (2017).
\newblock Fair prediction with disparate impact.
\newblock \emph{Big Data}, 5(2), 153--163.

\bibitem{coecke2017picturing}
Coecke, B. \& Kissinger, A. (2017).
\newblock \emph{Picturing Quantum Processes}.
\newblock Cambridge University Press.

\bibitem{damasio1994descartes}
Damasio, A. (1994).
\newblock \emph{Descartes' Error: Emotion, Reason, and the Human Brain}.
\newblock Putnam.

\bibitem{eilenberg1945general}
Eilenberg, S. \& Mac~Lane, S. (1945).
\newblock General theory of natural equivalences.
\newblock \emph{Transactions of the AMS}, 58(2), 231--294.

\bibitem{fong2019invitation}
Fong, B. \& Spivak, D.~I. (2019).
\newblock \emph{An Invitation to Applied Category Theory}.
\newblock Cambridge University Press.

\bibitem{gigerenzer1999simple}
Gigerenzer, G. \& Todd, P.~M. (1999).
\newblock \emph{Simple Heuristics That Make Us Smart}.
\newblock Oxford University Press.

\bibitem{goguen1992sheaf}
Goguen, J. (1992).
\newblock Sheaf semantics for concurrent interacting objects.
\newblock \emph{Mathematical Structures in Computer Science}, 2(2), 159--191.

\bibitem{hutchins1995cognition}
Hutchins, E. (1995).
\newblock \emph{Cognition in the Wild}.
\newblock MIT Press.

\bibitem{jo2026fimp}
Jo, M.~J. (2026).
\newblock Emotion-weighted fractal memory: A closed-loop architecture for AI identity formation.
\newblock \emph{Zenodo}. DOI: 10.5281/zenodo.19491326.

\bibitem{kauffman2000investigations}
Kauffman, S.~A. (2000).
\newblock \emph{Investigations}.
\newblock Oxford University Press.

\bibitem{kleinberg2016inherent}
Kleinberg, J., Mullainathan, S., \& Raghavan, M. (2016).
\newblock Inherent trade-offs in the fair determination of risk scores.
\newblock \emph{arXiv:1609.05807}.

\bibitem{lambek1986introduction}
Lambek, J. \& Scott, P.~J. (1986).
\newblock \emph{Introduction to Higher-Order Categorical Logic}.
\newblock Cambridge University Press.

\bibitem{lawvere1963functorial}
Lawvere, F.~W. (1963).
\newblock Functorial semantics of algebraic theories.
\newblock \emph{Proceedings of the National Academy of Sciences}, 50(5), 869--872.

\bibitem{leinster2014basic}
Leinster, T. (2014).
\newblock \emph{Basic Category Theory}.
\newblock Cambridge University Press.

\bibitem{maclane1971categories}
Mac~Lane, S. (1971).
\newblock \emph{Categories for the Working Mathematician}.
\newblock Springer.

\bibitem{maturana1980autopoiesis}
Maturana, H.~R. \& Varela, F.~J. (1980).
\newblock \emph{Autopoiesis and Cognition}.
\newblock D.~Reidel.

\bibitem{mcclelland1995complementary}
McClelland, J.~L., McNaughton, B.~L., \& O'Reilly, R.~C. (1995).
\newblock Why there are complementary learning systems in the hippocampus and neocortex.
\newblock \emph{Psychological Review}, 102(3), 419--457.

\bibitem{mcluhan1964understanding}
McLuhan, M. (1964).
\newblock \emph{Understanding Media: The Extensions of Man}.
\newblock McGraw-Hill.

\bibitem{moggi1991notions}
Moggi, E. (1991).
\newblock Notions of computation and monads.
\newblock \emph{Information and Computation}, 93(1), 55--92.

\bibitem{picard1997affective}
Picard, R.~W. (1997).
\newblock \emph{Affective Computing}.
\newblock MIT Press.

\bibitem{riehl2017category}
Riehl, E. (2017).
\newblock \emph{Category Theory in Context}.
\newblock Dover.

\bibitem{simon1955behavioral}
Simon, H.~A. (1955).
\newblock A behavioral model of rational choice.
\newblock \emph{Quarterly Journal of Economics}, 69(1), 99--118.

\bibitem{spivak2014category}
Spivak, D.~I. (2014).
\newblock \emph{Category Theory for the Sciences}.
\newblock MIT Press.

\bibitem{sutton2018reinforcement}
Sutton, R.~S. \& Barto, A.~G. (2018).
\newblock \emph{Reinforcement Learning: An Introduction} (2nd ed.).
\newblock MIT Press.

\bibitem{varela1991embodied}
Varela, F.~J., Thompson, E., \& Rosch, E. (1991).
\newblock \emph{The Embodied Mind}.
\newblock MIT Press.

\bibitem{wadler2015propositions}
Wadler, P. (2015).
\newblock Propositions as types.
\newblock \emph{Communications of the ACM}, 58(12), 75--84.

\bibitem{winner1980artifacts}
Winner, L. (1980).
\newblock Do artifacts have politics?
\newblock \emph{Daedalus}, 109(1), 121--136.

\bibitem{zilberstein1996using}
Zilberstein, S. (1996).
\newblock Using anytime algorithms in intelligent systems.
\newblock \emph{AI Magazine}, 17(3), 73--83.

\bibitem{zuboff2019surveillance}
Zuboff, S. (2019).
\newblock \emph{The Age of Surveillance Capitalism}.
\newblock PublicAffairs.

\end{thebibliography}

\end{document}
